%-----------------------------------------------------------------------
\subsection{The View Update Problem (Computer Science)}
\label{sec:instance-view-update}
%-----------------------------------------------------------------------

The \emph{view update problem} asks: given a database view (a projection of
the base relation), how can updates to the view be translated back to the
base relation?  \citet{bancilhon1981update} characterized exactly when unique
translation is possible via the \emph{complement view} framework: a view
admits a unique update translator if and only if a complement view exists such
that the pair forms a lossless decomposition.  As a corollary, non-injective
views without complements have no unique translation---a structural observation
that our framework captures as a minimal witness of the Explanation
Impossibility.

\paragraph{The explanation system.}
Define $\mathcal{S}_{\mathrm{view}} = (\Params, \sH, Y, \mathsf{obs},
\mathsf{exp}, \incomp)$ as follows:
\begin{itemize}
  \item $\Params = \sH = \texttt{Row}$ (database rows with two Boolean columns
    \texttt{colA} and \texttt{colB});
  \item $Y = \texttt{Bool}$ (the view: the projection onto \texttt{colA});
  \item $\mathsf{obs} = \texttt{view}$, where $\texttt{view}(r) = r.\texttt{colA}$;
  \item $\mathsf{exp} = \mathrm{id}$ (the full row is the ``explanation'' of
    the view);
  \item $\incomp(r_1, r_2) \iff r_1 \ne r_2$.
\end{itemize}

\paragraph{The Rashomon property.}
The witnesses are $r_1 = (\mathtt{true}, \mathtt{true})$ and
$r_2 = (\mathtt{true}, \mathtt{false})$, which satisfy
$\texttt{view}(r_1) = \mathtt{true} = \texttt{view}(r_2)$ and
$r_1 \ne r_2$.  The Rashomon property is constructively witnessed with
zero axioms.

\paragraph{The impossibility.}
By Theorem~\ref{thm:universal}, no translation from views back to base
relations can simultaneously be faithful (the translated row matches the
actual base row), stable (all rows with the same view receive the same
translation), and decisive (the translation distinguishes every pair of
distinct base rows).  This captures the impossibility corollary of \citet{bancilhon1981update}'s
complement framework: when no complement view exists, the non-injective view
admits no unique translation.  Our Lean formalization provides a minimal
constructive witness of this structural phenomenon, not a full encoding of the
complement lattice theory.

\paragraph{Resolution: complement views.}
The classical resolution \citep{bancilhon1981update} is to use
\emph{complement views}: additional projections that, together with the
original view, uniquely determine the base row.  In our framework, the
complement view is the additional $G$-invariant information needed to
collapse the fiber to a single point.  When no complement is available, the
resolution is to report the entire fiber: all base rows consistent with the
view.  This sacrifices decisiveness (no single row is selected) to achieve
stability and faithfulness.

\paragraph{Empirical illustration.}
Census disaggregation using real U.S.\ state and county structure (2020 Census
Bureau county counts, Zipf-distributed populations scaled to actual state
totals): the KL-divergence between the true county distribution and
Dirichlet-sampled alternatives scales strongly with the number of counties
(Spearman $\rho = 0.98$, $p = 7.9 \times 10^{-37}$; Pearson $r = 0.82$,
$p = 2.5 \times 10^{-13}$).  KL ranges from 0.00 (DC, 1 county) to 1.76
(Texas, 254 counties).  The negative control (DC) has KL~$= 0$ exactly.

\paragraph{Historical note.}
The view update problem was identified by \citet{bancilhon1981update} in 1981.
Our Lean formalization captures a minimal witness (a 2-row, 2-column Boolean
database with a single-column projection) that exhibits the core structural
phenomenon---non-injective observation implies Rashomon---rather than
formalizing the full complement lattice theory.  The structural parallel with
gauge invariance in physics, syntactic ambiguity in linguistics, and the phase
problem in crystallography was not previously recognized.

\paragraph{Lean formalization.}
The Lean~4 types are \texttt{Row} (a structure with fields
\texttt{colA colB : Bool}), \texttt{view} (the observation map
$r \mapsto r.\texttt{colA}$), and
\texttt{viewUpdateSystem : ExplanationSystem Row Row Bool}.
The impossibility is \texttt{view\_update\_impossibility} in
\texttt{ViewUpdate.lean}, derived with zero axioms.
